\documentclass[a4paper,12pt]{article}
\usepackage{polyglossia}
\setmainlanguage{arabic}
\setotherlanguages{italian,english}
\newfontfamily\arabicfont[Script=Arabic]{Amiri}
\usepackage{bidi}
\usepackage{geometry}
\usepackage{hyperref}
\geometry{margin=2.5cm}

\title{\textbf{سياسة الخصوصية المتعلقة بمعالجة البيانات الشخصية}\\
بموجب المادتين 13 و 14 من اللائحة (الاتحاد الأوروبي) 2016/679 (GDPR)\\
NutriApp}
\date{\today}

\begin{document}

    \maketitle

    \section*{1. مراقب البيانات}

    مراقب البيانات لمعالجة البيانات الشخصية هو \textbf{فريق NutriApp}، ومقره التشغيلي في إيطاليا.
    للحصول على أي معلومات تتعلق بمعالجة البيانات الشخصية، يمكنك الاتصال بمراقب البيانات على عنوان البريد الإلكتروني التالي:
    \href{mailto:privacy@nutriapp.it}{privacy@nutriapp.it}.

    إذا تم تعيين مسؤول حماية البيانات (DPO)، فسيتم نشر تفاصيل الاتصال الخاصة به وإبلاغ المستخدمين بها.

    \section*{2. نوع البيانات المعالجة}

    كجزء من استخدام تطبيق NutriApp، قد تتم معالجة الفئات التالية من البيانات الشخصية:

    \textbf{بيانات الهوية والاتصال}
    الاسم، عنوان البريد الإلكتروني، أي بيانات اعتماد تسجيل الدخول.

    \textbf{البيانات المتعلقة بالصحة (فئات خاصة من البيانات بموجب المادة 9 من GDPR)}
    الوزن، الطول، الأهداف الغذائية، عادات الأكل، البيانات المتعلقة بالسعرات الحرارية والمدخول الغذائي، أي معلومات يقدمها المستخدم طواعية فيما يتعلق بالحالات الطبية أو الاحتياجات الغذائية المحددة.

    \textbf{بيانات الاستخدام}
    المعلومات المتعلقة باستخدام التطبيق، التفضيلات، التفاعلات مع الميزات.

    \textbf{البيانات الفنية}
    عنوان IP، معرفات الجهاز، نظام التشغيل، بيانات السجل، ملفات تعريف الارتباط الفنية حيثما ينطبق ذلك.

    \section*{3. الغرض من المعالجة}

    تتم معالجة البيانات الشخصية للأغراض التالية:

    \begin{itemize}
        \item توفير خدمات المراقبة الغذائية وإدارة يوميات الطعام؛
        \item السماح بإنشاء وإدارة حساب المستخدم؛
        \item معالجة الإحصائيات المخصصة والتصورات الرسومية للتقدم؛
        \item تحسين جودة ووظائف التطبيق؛
        \item الوفاء بالالتزامات القانونية أو طلبات السلطات المختصة؛
        \item إدارة أي طلبات للحصول على المساعدة أو الدعم الفني.
    \end{itemize}

    \section*{4. الأساس القانوني للمعالجة}

    تستند معالجة البيانات الشخصية إلى الأسس القانونية التالية:

    \begin{itemize}
        \item المادة 6، الفقرة 1، الحرف (ب) من GDPR: تنفيذ عقد أو تدابير ما قبل التعاقد؛
        \item المادة 6، الفقرة 1، الحرف (ج) من GDPR: الامتثال للالتزامات القانونية؛
        \item المادة 6، الفقرة 1، الحرف (أ) من GDPR: موافقة صاحب البيانات لأغراض محددة؛
        \item المادة 9، الفقرة 2، الحرف (أ) من GDPR: الموافقة الصريحة على معالجة البيانات المتعلقة بالصحة.
    \end{itemize}

    يعد توفير البيانات ضروريًا لتوفير الخدمات الرئيسية للتطبيق. قد يؤدي الفشل في توفير البيانات إلى استحالة استخدام ميزات معينة.

    \section*{5. طرق المعالجة}

    تتم معالجة البيانات باستخدام تكنولوجيا المعلومات وأدوات الاتصال عن بعد، وفقًا لمبادئ الشرعية والإنصاف والشفافية وتقليل البيانات وتقييد التخزين.

    لا يتم تنفيذ أي عمليات صنع قرار آلية تنتج آثارًا قانونية على صاحب البيانات بموجب المادة 22 من GDPR.

    \section*{6. فترة الاحتفاظ بالبيانات}

    سيتم الاحتفاظ بالبيانات الشخصية للوقت الضروري للغاية لتحقيق الأغراض التي تم جمعها من أجلها، امتثالاً لمبدأ تقييد التخزين المشار إليه في المادة 5 من GDPR.

    بالأخص:

    \begin{itemize}
        \item سيتم الاحتفاظ بالبيانات طوال مدة العلاقة التعاقدية واستخدام التطبيق؛

        \item بمجرد إنهاء العلاقة التعاقدية أو حذف الحساب، يجوز الاحتفاظ بالبيانات الشخصية لفترة إضافية إذا كان ذلك ضروريًا من أجل:
        \begin{itemize}
            \item الوفاء بالالتزامات القانونية (على سبيل المثال، الالتزامات الضريبية أو المحاسبية)؛
            \item حماية حقوق مراقب البيانات، بما في ذلك في المحكمة؛
            \item إدارة أي شكاوى أو نزاعات.
        \end{itemize}

        \item سيتم الاحتفاظ بالبيانات المتعلقة بالصحة حصريًا للوقت اللازم لتقديم الخدمة التي طلبها المستخدم وسيتم حذفها بشكل آمن أو جعلها مجهولة المصدر عند إنهاء العلاقة، ما لم يكن هناك التزام قانوني أو أساس قانوني آخر يسمح بالاحتفاظ بها.
    \end{itemize}

    في نهاية الفترات المشار إليها أعلاه، سيتم حذف البيانات بشكل آمن أو جعلها مجهولة المصدر بشكل لا رجعة فيه.

    \section*{7. نقل البيانات وتوصيلها}

    يجوز إبلاغ البيانات الشخصية إلى:

    \begin{itemize}
        \item مزودي خدمات تكنولوجيا المعلومات والاستضافة؛
        \item مزودي البنية التحتية السحابية؛
        \item المستشارين القانونيين أو الضريبيين، عند الضرورة.
    \end{itemize}

    تعمل هذه الكيانات كمعالجين للبيانات بموجب المادة 28 من GDPR.

    إذا تم نقل البيانات إلى دول ثالثة خارج المنطقة الاقتصادية الأوروبية، فسيتم النقل وفقًا للمادتين 44 وما يليها من GDPR، من خلال اعتماد الضمانات المناسبة (على سبيل المثال، البنود التعاقدية القياسية).

    \section*{8. حقوق صاحب البيانات}

    يجوز لصاحب البيانات ممارسة الحقوق المنصوص عليها في المواد 15-22 من GDPR في أي وقت، بما في ذلك:

    \begin{itemize}
        \item حق الوصول إلى البيانات الشخصية؛
        \item الحق في التصحيح؛
        \item الحق في المحو (الحق في النسيان)؛
        \item الحق في تقييد المعالجة؛
        \item الحق في قابلية نقل البيانات؛
        \item الحق في الاعتراض على المعالجة؛
        \item الحق في سحب الموافقة في أي وقت.
    \end{itemize}

    يمكن إرسال الطلبات إلى عنوان البريد الإلكتروني \href{mailto:privacy@nutriapp.it}{privacy@nutriapp.it}.

    من الممكن أيضًا تقديم شكوى إلى هيئة حماية البيانات.

    \section*{9. أمن البيانات}

    يعتمد NutriApp التدابير الفنية والتنظيمية المناسبة بموجب المادة 32 من GDPR، بما في ذلك:

    \begin{itemize}
        \item المصادقة الآمنة للمستخدم؛
        \item يقتصر الوصول إلى البيانات على الموظفين المصرح لهم؛
        \item عمليات النسخ الاحتياطي الدورية وأنظمة الحماية ضد الوصول غير المصرح به.
    \end{itemize}

    \section*{10. القصر}

    التطبيق غير مخصص للأطفال دون سن 14 عامًا.

    بموجب المادة 8 من GDPR والقوانين المحلية المعمول بها، فإن معالجة البيانات الشخصية بناءً على الموافقة فيما يتعلق بالعرض المباشر لخدمات مجتمع المعلومات تكون قانونية فقط إذا كان عمر الطفل 14 عامًا على الأقل.

    إذا كان المستخدم أقل من 14 عامًا، فإن معالجة البيانات الشخصية تكون قانونية فقط إذا تم منح الموافقة أو الإذن بها من قبل صاحب المسؤولية الأبوية.

    يجوز لمراقب البيانات اعتماد تدابير معقولة للتحقق من العمر الذي أعلنه المستخدم، وحيثما كان ذلك ضروريًا، التحقق من وجود موافقة من الوالد أو الوصي.

    إذا تبين أنه تم جمع بيانات شخصية لقاصر أقل من 14 عامًا دون الموافقة المطلوبة، فسيتم حذف هذه البيانات على الفور.

    \section*{11. التغييرات في سياسة الخصوصية هذه}

    قد يتم تحديث سياسة الخصوصية هذه بمرور الوقت.
    سيتم إبلاغ المستخدمين بأي تغييرات جوهرية عبر إشعار داخل التطبيق أو البريد الإلكتروني.

    \vspace{1cm}

    \noindent
    آخر تحديث: \today

\end{document}
